% Standalone section. Packages: amsmath,amssymb,amsthm,array,longtable,booktabs.
% A minimal compiling wrapper is wrapper.tex in this directory.
\newtheorem{htthm}{Theorem}[section]
\newtheorem{htlem}[htthm]{Lemma}
\newtheorem{htdef}[htthm]{Definition}
\newtheorem{htprop}[htthm]{Proposition}
\newtheorem{htrem}[htthm]{Remark}
\newcommand{\htU}{\mathrm U}
\newcommand{\htD}{\mathrm D}
\newcommand{\htX}{\mathrm X}
\newcommand{\htM}[2]{#1_{\rm m}(#2)}
\newcommand{\htV}[2]{#1[#2]}
\newcommand{\htcur}{Z_{\rm cur}}
\newcommand{\htnext}{Z_{\rm next}}
\newcommand{\htF}{F_c}
\newcommand{\htball}{\overline B}
\section{A tree input for the actual local passage}\label{ht:section}

References prefixed by P refer to the supplied version of the paper
\emph{The spectral gap of the spin-2 AKLT model on the square lattice}.
In particular, P.Definitions 2.1, 2.3, 6.1, 6.2, 6.8 and 6.9 define,
respectively, the laws, pins, domains, trace, covered patterns and (H1).
P.Lemmas 7.2 and 7.3 identify the non-exceptional departures.
We prove the inequality needed at those departures, with
\begin{equation}\label{ht:constants}
 c_{\rm s}=961/1000,\qquad c_{\rm t}=187/200.
\end{equation}
The replacement is a computer-assisted statement with explicitly specified
finite rational checks. It does not require the old pattern catalogue.

\begin{htrem}[Scope of the replacement]\label{ht:scope}
P.Definition 6.8 fills two external slots at every D corner, even when
that corner has two present internal bonds. An actual D corner cannot
have these additional factors: its kept set has size two and it carries
no original pin. Consequently the argument below proves the actual
local conclusion of P.Lemma 6.10, not the larger covered-pattern assertion
in P.Definition 6.9. In uses of (H1) in Sections 6--7, replace its
hypothesis by Theorem~\ref{ht:reduction}. P.Lemmas 7.2 and 7.3 supply
its geometric hypotheses, and the same exceptional bound is retained.
No assertion that the larger covered-pattern inequality is false is made.
\end{htrem}

\begin{htdef}[Laws and local factors]\label{ht:laws}
Let \(D\) be the twelve unit vectors of P.Definition 2.1 and let
\(\lambda_{\htU}\) be the uniform probability on \(D\). Put
\(\lambda_{\htD}=(3/5)\lambda_{\htU}+(2/5)\delta_0\).
Integration always means a finite sum. A field \(a\), \(|a|\le1\),
contributes \(1+a\cdot n\), and a bond contributes \(1+n\cdot n'\).
The moments needed here, from P.Lemma 2.2, are
\begin{equation}\label{ht:moments}
 \lambda_{\htU}(n_i)=\lambda_{\htU}(n_in_jn_k)=0,\quad
 \lambda_{\htU}(n_in_j)=\delta_{ij}/3,\quad
 \lambda_{\htU}(n_in_jn_kn_l)
 =\frac{\delta_{ij}\delta_{kl}+\delta_{ik}\delta_{jl}
                +\delta_{il}\delta_{jk}}{15}.
\end{equation}
Both laws have constant moment one. Every positive-degree D moment
through degree four is \(3/5\) times the corresponding U moment.
These tensors, including their orthogonal covariance, can be used at
each site whenever the separate degree is at most four.
\end{htdef}

\begin{htlem}[Positive normalizers, affinity, and leaf summation]
\label{ht:elementary}
A product of at most four field or bond factors at a site has a strictly
positive single-site normalizer when the other spins are fixed in the
closed unit ball. The same holds for a finite network with at most four
incident factors at every site. If a leaf has law \(\lambda\), external
weight \(W\), normalizer \(A=\lambda(W)\), and mean
\(\mu=\lambda(nW)/A\), integrating its bond to \(x\) gives
\begin{equation}\label{ht:leaf}
 A(1+\mu\cdot n_x).
\end{equation}
With at most two unit-ball fields a U mean has norm at most \(1/2\).
With at most one unit-ball field a D mean has norm at most \(1/5\);
with its sole field of radius \(3/5\) it has norm at most \(3/25\).
With at most three unit-ball fields either law has mean norm at most
\(3/5\).

An inequality separately affine in vectors and valid when their norms
are the listed radii is valid throughout their product of closed balls.
Inserting a zero field is exact. Enlarging a listed radius is legitimate
in this sense, without any monotonicity assertion about partition ratios.
\end{htlem}
\begin{proof}
All factors are nonnegative. A factor with a vector of norm at most one
can exclude at most one of the twelve U letters. Greedily choosing a
letter at each site excludes at most four letters and produces a strictly
positive summand, exactly as in P.Lemma 2.4. This also proves the
single-site assertion and covers D laws because they give positive mass
to every U letter. Formula~\eqref{ht:leaf} is linearity of integration.
For two unit U fields \(a,b\), write \(t=a\cdot b\). The normalizer
and unnormalized mean are \(A=1+t/3\), \(V=(a+b)/3\), and
\(A^2-4|V|^2=(1-t)^2/9\ge0\). For one D field \(a\), they are
\(A=1\), \(V=a/5\). The three-field bound is P.Lemma 2.5.

For radius \(r>0\), any \(|a|\le r\) has, for some unit \(e\),
\begin{equation}\label{ht:convexfield}
 1+a\cdot n=\frac{1+|a|/r}{2}(1+re\cdot n)
             +\frac{1-|a|/r}{2}(1-re\cdot n).
\end{equation}
Both coefficients are nonnegative and sum to one. When \(a=0\),
choose any \(e\); when \(r=0\) there is nothing to prove.
Expansion proves the affine assertion and also fills missing factors.
For a normalized mean this expansion gives weights
\(t_\eta A_\eta/\sum_\xi t_\xi A_\xi\), all nonnegative.
Thus bounds on means extend to balls as well. Normalizers are positive
by the first part. The displayed two-field computation therefore covers
fewer fields and all field norms at most one.
\end{proof}

\begin{htdef}[Tree models and their passage polynomials]\label{ht:models}
A model is a disjoint union of paths. A spin vertex has law U or D,
a list of field radii in \([0,1]\), and a mark indicating that it was a
neighbour of the passage site. At most three vertices are marked.
At most one of them has an outgoing mark, which is allowed to be absent.
An isolated marked mean vertex has instead a deterministic vector
\(m\in\htball(0,R)\), with its specified radius \(R\), and contributes
no integration or field factor. A bond joins consecutive spin vertices;
mean vertices are isolated. Each listed radius \(r\) denotes the factor
\(1+ru\cdot n\), where its unit direction \(u\) is arbitrary.

Let \(\mathcal W\) be the product of the laws, path bonds and field
factors, and write
\begin{equation}\label{ht:polynomials}
 Z=\int\mathcal W,\qquad
 C_{ij}=\int (n_i\cdot n_j)\mathcal W,\qquad
 \htcur=Z+\frac13\sum_{\{i,j\}\subset M}C_{ij},\qquad
 \htnext=Z+\frac15 C_{fg}.
\end{equation}
Here \(M\) is the marked set, insertion at a mean vertex means insertion
of its deterministic vector, and the last term is present exactly when
there are two non-outgoing marked vertices \(f,g\). If the outgoing mark
is absent, all marked vertices are non-outgoing; our list then has at
most two marks. In every case without exactly two non-outgoing marks, the last term is
zero. Put \(F_c=\htcur-c\htnext\). All listed mean radii are at
most one. Thus \(Z>0\), \(\htnext\ge4Z/5>0\), and
\(\htcur\ge Z/2>0\): the latter follows from
\(\sum_{i<j}n_i\cdot n_j=(|\sum_i n_i|^2-\sum_i|n_i|^2)/2
\ge-3/2\).

The precise list is Table~\ref{ht:modeltable}, with zero-based indices.
Coordinates are only convenient names:
\[
 \begin{gathered}
 E=(1,0),\ A=(1,1),\ N=(0,1),\ B=(-1,1),\ W=(-1,0),\\
 L=(-1,-1),\ S=(0,-1),\ R=(1,-1).
 \end{gathered}
\]
Every included \(E,N,W\) is marked. The outgoing mark is \(E\) in a
turning row and \(N\) in a straight row, if that vertex is included.
\(\htM{x}{r}\) means an isolated mean vertex of radius \(r\).
\(\htV{x}{r_1,\ldots,r_k}\) means a U spin vertex with those radii.
The undecorated letters \(A,B,L,R\) mean unmarked U vertices with
radii \((1,1)\), and \(S\) means an unmarked D vertex without fields.
Hyphens join consecutive vertices; semicolons separate components.
Zero radii are displayed and are literal zero fields.
\end{htdef}

\begin{htdef}[H1-tree]\label{ht:input}
For each of the 38 rows of Table~\ref{ht:modeltable}, for every independent
choice of all unit field directions and all mean vectors in their
listed balls, require
\begin{equation}\label{ht:h1tree}
 \htcur\ge c_{\rm type}\htnext.
\end{equation}
Equivalently one may require each mean vector on the boundary of its
ball, by Lemma~\ref{ht:elementary}. This is the finite list of universal
inequalities called H1-tree.
\end{htdef}

\begingroup\small
\begin{longtable}{@{}rcp{0.84\textwidth}@{}}
\caption{The complete finite input.}\label{ht:modeltable}\\
\toprule Index & Type & Components\\\midrule\endfirsthead
\toprule Index & Type & Components (continued)\\\midrule\endhead
\bottomrule\endfoot
0 & t & \(\htM{E}{\frac{11329}{25630}}\)\\[2pt]
1 & t & \(\htM{W}{\frac{167}{440}}\ ;\quad \htM{E}{\frac{226}{495}}\)\\[2pt]
2 & t & \(\htM{W}{\frac{226}{495}}\ ;\quad \htM{E}{\frac{167}{440}}\)\\[2pt]
3 & t & \(\htM{W}{\frac{169}{385}}\ ;\quad \htM{E}{\frac{169}{385}}\)\\[2pt]
4 & t & \(\htM{W}{\frac{11329}{25630}}\)\\[2pt]
5 & t & \(\htV{N}{\frac{1}{2},\frac{3}{5}}\! -\! A\! -\! \htV{E}{\frac{137}{265},\frac{3}{5}}\)\\[2pt]
6 & t & \(\htM{N}{\frac{10}{39}}\ ;\quad \htM{W}{\frac{5}{14}}\ ;\quad \htM{E}{\frac{1}{5}}\)\\[2pt]
7 & t & \(\htM{N}{\frac{1}{5}}\ ;\quad \htM{W}{\frac{5}{14}}\ ;\quad \htM{E}{\frac{5}{14}}\)\\[2pt]
8 & t & \(\htM{N}{\frac{3}{25}}\ ;\quad \htM{W}{\frac{451}{1125}}\ ;\quad \htM{E}{\frac{451}{1125}}\)\\[2pt]
9 & t & \(\htM{N}{\frac{1}{6}}\ ;\quad \htM{W}{\frac{22429}{58525}}\)\\[2pt]
10 & t & \(\htM{W}{\frac{137}{795}}\ ;\quad \htM{N}{\frac{167}{440}}\)\\[2pt]
11 & t & \(\htM{N}{\frac{1}{3}}\ ;\quad \htM{W}{\frac{740}{2193}}\)\\[2pt]
12 & t & \(\htM{W}{\frac{451}{1125}}\ ;\quad \htV{N}{0,0}\! -\! A\! -\! \htV{E}{\frac{3}{5},\frac{3}{5}}\)\\[2pt]
13 & t & \(\htM{W}{\frac{1}{5}}\ ;\quad \htV{N}{\frac{1}{5},\frac{3}{5}}\! -\! A\! -\! \htV{E}{\frac{3}{5},\frac{3}{5}}\)\\[2pt]
14 & t & \(\htM{W}{\frac{5}{14}}\ ;\quad \htV{N}{0,\frac{3}{5}}\! -\! A\! -\! \htV{E}{\frac{3}{5},\frac{3}{5}}\)\\[2pt]
15 & t & \(\htV{N}{\frac{1}{2},\frac{3}{5}}\! -\! B\! -\! \htV{W}{\frac{137}{265},\frac{3}{5}}\)\\[2pt]
16 & t & \(\htM{E}{\frac{1}{5}}\ ;\quad \htV{N}{\frac{1}{5},\frac{3}{5}}\! -\! B\! -\! \htV{W}{\frac{3}{5},\frac{3}{5}}\)\\[2pt]
17 & t & \(\htM{E}{\frac{451}{1125}}\ ;\quad \htV{N}{0,0}\! -\! B\! -\! \htV{W}{\frac{3}{5},\frac{3}{5}}\)\\[2pt]
18 & t & \(\htM{E}{\frac{5}{14}}\ ;\quad \htV{N}{0,\frac{3}{5}}\! -\! B\! -\! \htV{W}{\frac{3}{5},\frac{3}{5}}\)\\[2pt]
19 & t & \(\htV{W}{\frac{3}{5},\frac{3}{5}}\! -\! B\! -\! \htV{N}{\frac{3}{5}}\! -\! A\! -\! \htV{E}{\frac{3}{5},\frac{3}{5}}\)\\[2pt]
20 & t & \(\htV{W}{\frac{1}{2},\frac{3}{5}}\! -\! L\! -\! S\! -\! R\! -\! \htV{E}{\frac{1}{2},\frac{3}{5}}\)\\[2pt]
21 & s & \(\htM{E}{\frac{11329}{25630}}\)\\[2pt]
22 & s & \(\htM{W}{\frac{137}{795}}\ ;\quad \htM{N}{\frac{167}{440}}\)\\[2pt]
23 & s & \(\htM{E}{\frac{22429}{58525}}\ ;\quad \htM{N}{\frac{1}{6}}\)\\[2pt]
24 & s & \(\htM{E}{\frac{740}{2193}}\ ;\quad \htM{N}{\frac{1}{3}}\)\\[2pt]
25 & s & \(\htM{N}{\frac{169}{385}}\)\\[2pt]
26 & s & \(\htV{E}{\frac{137}{265},\frac{3}{5}}\! -\! A\! -\! \htV{N}{\frac{1}{2},\frac{3}{5}}\)\\[2pt]
27 & s & \(\htM{W}{\frac{167}{440}}\ ;\quad \htM{E}{\frac{169}{385}}\)\\[2pt]
28 & s & \(\htM{W}{\frac{317}{1025}}\ ;\quad \htM{E}{\frac{226}{495}}\)\\[2pt]
29 & s & \(\htM{E}{\frac{1}{5}}\ ;\quad \htM{W}{\frac{1}{5}}\ ;\quad \htM{N}{\frac{317}{1025}}\)\\[2pt]
30 & s & \(\htM{W}{\frac{1}{5}}\ ;\quad \htM{E}{\frac{5}{14}}\ ;\quad \htM{N}{\frac{10}{39}}\)\\[2pt]
31 & s & \(\htM{E}{\frac{451}{1125}}\ ;\quad \htM{W}{\frac{451}{1125}}\ ;\quad \htM{N}{\frac{3}{25}}\)\\[2pt]
32 & s & \(\htM{E}{\frac{5}{14}}\ ;\quad \htM{W}{\frac{5}{14}}\ ;\quad \htM{N}{\frac{1}{5}}\)\\[2pt]
33 & s & \(\htM{W}{\frac{1}{5}}\ ;\quad \htV{E}{\frac{3}{5},\frac{3}{5}}\! -\! A\! -\! \htV{N}{\frac{1}{5},\frac{3}{5}}\)\\[2pt]
34 & s & \(\htM{E}{\frac{5}{14}}\ ;\quad \htV{W}{\frac{3}{5},\frac{3}{5}}\! -\! B\! -\! \htV{N}{0,\frac{3}{5}}\)\\[2pt]
35 & s & \(\htM{W}{\frac{451}{1125}}\ ;\quad \htV{E}{\frac{3}{5},\frac{3}{5}}\! -\! A\! -\! \htV{N}{0,0}\)\\[2pt]
36 & s & \(\htV{W}{\frac{3}{5},\frac{3}{5}}\! -\! B\! -\! \htV{N}{\frac{3}{5}}\! -\! A\! -\! \htV{E}{\frac{3}{5},\frac{3}{5}}\)\\[2pt]
37 & s & \(\htV{W}{\frac{1}{5},\frac{3}{5}}\! -\! L\! -\! S\! -\! R\! -\! \htV{E}{\frac{1}{2},\frac{3}{5}}\)\\[2pt]
\end{longtable}
\endgroup

\begin{htlem}[Finite unequal-radius mean bounds]\label{ht:meanbern}
For each row of Table~\ref{ht:meantable}, a site with its indicated law,
three fields of norms at most \((r_1,r_2,r_3)\), and no other factors,
has normalized mean of norm at most the listed \(R\). A zero radius
permits an absent factor. These 21 statements follow from 567 explicit
nonnegative rational Bernstein coefficients, specified below.
\end{htlem}
\begin{proof}
First put all nonzero fields on their prescribed spheres, and denote
them by \(a,b,c\). Set
\[
 x=a\cdot b,\quad y=a\cdot c,\quad z=b\cdot c,\quad
 V=(5+z)a+(5+y)b+(5+x)c,\quad
 d=\begin{cases}3&\htU,\\5&\htD.\end{cases}
\]
Equation~\eqref{ht:moments} gives exactly
\begin{equation}\label{ht:meanformula}
 \mu=\frac{V}{5(d+x+y+z)},\qquad
 R=\frac{(r_1+r_2+r_3)/3+r_1r_2r_3/5}
 {d/3+(r_1r_2+r_1r_3+r_2r_3)/3}.
\end{equation}
The denominator of the mean is positive by Lemma~\ref{ht:elementary}.
Consider
\begin{align}\label{ht:bernpoly}
 P(x,y,z)={}&25R^2(d+x+y+z)^2
  -(5+z)^2r_1^2-(5+y)^2r_2^2-(5+x)^2r_3^2\notag\\
 &-2x(5+z)(5+y)-2y(5+z)(5+x)-2z(5+y)(5+x).
\end{align}
Make the substitutions
\(x=r_1r_2(1-2s_1)\), \(y=r_1r_3(1-2s_2)\),
\(z=r_2r_3(1-2s_3)\). The resulting polynomial has degree at most two
in each \(s_i\). If its monomial coefficients are \(p_\alpha\), its
Bernstein coefficients of degree \((2,2,2)\) are exactly
\begin{equation}\label{ht:bernformula}
 b_\beta=\sum_{\alpha\le\beta}p_\alpha
           \prod_{j=1}^3\frac{\binom{\beta_j}{\alpha_j}}
                                 {\binom2{\alpha_j}},
 \qquad \beta\in\{0,1,2\}^3.
\end{equation}
The finite assertion is that each of these 27 coefficients is
nonnegative for each of the 21 rows. The file
\texttt{data/\allowbreak universal\_mean\_checks.json} lists all 567 coefficients;
\texttt{checks/audit.py} reconstructs \eqref{ht:bernpoly} and
\eqref{ht:bernformula}, compares every coefficient with that file, and
checks its sign using rational arithmetic. The largest numerator or
denominator has 53 bits. Nonnegative Bernstein coefficients give
\(P\ge0\) on the full cube, and hence on every actual Gram triple,
without needing a Gram determinant constraint. Taking square roots
in \eqref{ht:meanformula} proves the bound on spheres. The last paragraph
of Lemma~\ref{ht:elementary} extends it to balls. We do not assert
\eqref{ht:meanformula}'s proposed radius for arbitrary triples of radii.
\end{proof}
\begin{longtable}{@{}cccccc@{}}
\caption{All 21 unequal-radius checks.}\label{ht:meantable}\\
\toprule Law & \(r_1\)&\(r_2\)&\(r_3\)&\(R\)& Bernstein degree\\\midrule\endfirsthead
\toprule Law & \(r_1\)&\(r_2\)&\(r_3\)&\(R\)& Bernstein degree\\\midrule\endhead
\bottomrule\endfoot
\(\htD\) & \(0\) & \(0\) & \(\frac{3}{5}\) & \(\frac{3}{25}\) & \((2,2,2)\)\\
\(\htU\) & \(0\) & \(0\) & \(0\) & \(0\) & \((2,2,2)\)\\
\(\htU\) & \(0\) & \(0\) & \(\frac{1}{5}\) & \(\frac{1}{15}\) & \((2,2,2)\)\\
\(\htU\) & \(0\) & \(0\) & \(\frac{1}{2}\) & \(\frac{1}{6}\) & \((2,2,2)\)\\
\(\htU\) & \(0\) & \(0\) & \(\frac{18}{35}\) & \(\frac{6}{35}\) & \((2,2,2)\)\\
\(\htU\) & \(0\) & \(0\) & \(\frac{137}{265}\) & \(\frac{137}{795}\) & \((2,2,2)\)\\
\(\htU\) & \(0\) & \(0\) & \(\frac{3}{5}\) & \(\frac{1}{5}\) & \((2,2,2)\)\\
\(\htU\) & \(0\) & \(\frac{1}{5}\) & \(\frac{3}{5}\) & \(\frac{10}{39}\) & \((2,2,2)\)\\
\(\htU\) & \(0\) & \(\frac{1}{2}\) & \(\frac{3}{5}\) & \(\frac{1}{3}\) & \((2,2,2)\)\\
\(\htU\) & \(0\) & \(\frac{18}{35}\) & \(\frac{3}{5}\) & \(\frac{65}{193}\) & \((2,2,2)\)\\
\(\htU\) & \(0\) & \(\frac{137}{265}\) & \(\frac{3}{5}\) & \(\frac{740}{2193}\) & \((2,2,2)\)\\
\(\htU\) & \(0\) & \(\frac{3}{5}\) & \(\frac{3}{5}\) & \(\frac{5}{14}\) & \((2,2,2)\)\\
\(\htU\) & \(\frac{1}{5}\) & \(\frac{1}{5}\) & \(\frac{3}{5}\) & \(\frac{317}{1025}\) & \((2,2,2)\)\\
\(\htU\) & \(\frac{1}{5}\) & \(\frac{1}{2}\) & \(\frac{3}{5}\) & \(\frac{167}{440}\) & \((2,2,2)\)\\
\(\htU\) & \(\frac{1}{5}\) & \(\frac{18}{35}\) & \(\frac{3}{5}\) & \(\frac{2956}{7725}\) & \((2,2,2)\)\\
\(\htU\) & \(\frac{1}{5}\) & \(\frac{137}{265}\) & \(\frac{3}{5}\) & \(\frac{22429}{58525}\) & \((2,2,2)\)\\
\(\htU\) & \(\frac{1}{5}\) & \(\frac{3}{5}\) & \(\frac{3}{5}\) & \(\frac{451}{1125}\) & \((2,2,2)\)\\
\(\htU\) & \(\frac{1}{2}\) & \(\frac{1}{2}\) & \(\frac{3}{5}\) & \(\frac{169}{385}\) & \((2,2,2)\)\\
\(\htU\) & \(\frac{1}{2}\) & \(\frac{18}{35}\) & \(\frac{3}{5}\) & \(\frac{2987}{6765}\) & \((2,2,2)\)\\
\(\htU\) & \(\frac{1}{2}\) & \(\frac{137}{265}\) & \(\frac{3}{5}\) & \(\frac{11329}{25630}\) & \((2,2,2)\)\\
\(\htU\) & \(\frac{1}{2}\) & \(\frac{3}{5}\) & \(\frac{3}{5}\) & \(\frac{226}{495}\) & \((2,2,2)\)\\
\end{longtable}

\begin{htlem}[Two unit fields and one small field]\label{ht:specialmean}
A U site with two fields of norm at most one and a third of norm at most
\(r\) has mean norm at most \((5+4r)/(10+5r)\) when
\begin{equation}\label{ht:specialradii}
 (r,R)=(1/10,18/35)\quad\hbox{or}\quad(3/25,137/265).
\end{equation}
Each assertion has a 25-coefficient rational Bernstein proof of degree
\((4,4)\).
\end{htlem}
\begin{proof}
By positive expansion it is enough that the first two fields \(b,c\)
are unit. Write \(t=|b+c|/2\in[0,1]\). Choose orthonormal directions
along \(b+c\) and \(b-c\) when nonzero, completing arbitrarily at
the endpoints. For the third field write \(a=\alpha u+\beta v+\gamma w\),
so \(\alpha^2+\beta^2+\gamma^2\le r^2\). Put \(d=2+2t^2\).
The vector \(V\) in \eqref{ht:meanformula} has these components:
\[
 2d\alpha+10t,\qquad (2+4t^2)\beta,\qquad (d+2)\gamma.
\]
Since \(2+4t^2\le d+2\),
\begin{equation}\label{ht:specialpoly}
 |V|^2\le(2d\alpha+10t)^2+(d+2)^2(r^2-\alpha^2),\qquad
 3Z=d+2t\alpha>0.
\end{equation}
Set \(t=s\), \(\alpha=-r+2rv\). The polynomial
\[
 25R^2(d+2t\alpha)^2-(2d\alpha+10t)^2
                         -(d+2)^2(r^2-\alpha^2)
\]
has degree at most four in each of \(s,v\). Apply
\eqref{ht:bernformula} with two variables and denominator binomial
\(\binom4{\alpha_j}\). Its 25 coefficients are nonnegative for each
pair in \eqref{ht:specialradii}; they are reconstructed and listed in
\texttt{data/mean\_receipt.json}. Their numerators and denominators need
at most 27 bits. Thus \(|\mu|=|V|/[5(d+2t\alpha)]\le R\).
The coordinate calculation and bound extend continuously to \(t=0,1\).
Positive expansion handles all smaller first and second fields.
The supplied \texttt{mean\_unit\_pair.json} also records \(r=1/25\),
which is checked by the replay but is not used here. Its original
three-variable storage repeats each two-variable coefficient five times.
\end{proof}

\begin{htlem}[The 22 scalar inequalities]\label{ht:scalar}
H1-tree holds for the 22 rows consisting entirely of mean vertices.
The finite certificate has 89 rational comparisons.
\end{htlem}
\begin{proof}
Give a missing outgoing vector radius \(r_o=0\); give a missing
non-outgoing vector radius zero. Permute the two non-outgoing means
as necessary. On their boundary spheres put
\[
 k=r_fr_g(1/3-c/5),\quad L(t)=1-c+kt,\quad
 Q(t)=L(t)^2-\frac{r_o^2}{9}
                  (r_f^2+r_g^2+2r_fr_gt),\quad -1\le t\le1.
\]
All components have base normalizer one. Minimizing the linear term
in the outgoing mean gives
\[
 F_c\ge L(t)-\frac{r_o}{3}
                    \sqrt{r_f^2+r_g^2+2r_fr_gt}.
\]
For each scalar row the finite tests are
\begin{equation}\label{ht:scalartests}
 L(-1),L(1),Q(-1),Q(1)\ge0,
 \quad Q(-q_1/(2q_2))\ge0
 \ \hbox{if }q_2>0\hbox{ and }-1<-q_1/(2q_2)<1,
\end{equation}
where \(Q=q_0+q_1t+q_2t^2\). Linear and quadratic minimization on
an interval proves \(L,Q\ge0\), hence \(F_c\ge0\). Boundary radii
zero cause no problem; affinity extends the conclusion to all balls.
The rows are indices
\(0,1,2,3,4,6,7,8,9,10,11,21,22,23,24,25,27,28,29,30,31,32\).
\texttt{data/\allowbreak universal\_certificates.jsonl} lists their radii,
\((q_0,q_1,q_2)\), and every tested value; the replay recomputes and
compares all of them. There are four tests per row and one interior
vertex test in total, making 89 comparisons. In particular the 19
older boxes in \texttt{closed\_boxes.json} are not the scalar
certificate for the present 38-row family.
\end{proof}

\begin{htdef}[Exact reconstruction of the other 16 inequalities]
\label{ht:reconstruct}
Omit zero fields. Order sites as \((E,A,N,B,W,L,S,R)\), and order the
remaining fields at each site as printed in Table~\ref{ht:modeltable}.
An isolated mean \(m=Ru\) can, for reconstruction only, be encoded as
a formal isolated U site with field \(3Ru\) and a mark:
\[
 \lambda_{\htU}(1+3Ru\cdot n)=1,\qquad
 \lambda_{\htU}(n(1+3Ru\cdot n))=Ru.
\]
The number \(3R\) may exceed one; it is an algebraic encoding, not a
physical-field assumption, and that site has degree at most two.

For an explicit coefficient recipe assign an index in \(\{0,1,2,3\}\)
to each path bond, one to each field, and a mark index to each marked
spin site. Index zero denotes the constant 1 and indices 1--3 denote
coordinates. At each site take the moment \eqref{ht:moments} of all
its nonzero indices. Multiply a field index zero by 1 and a positive
field index by its radius. Sum each bond index over 0--3. This contracts
\(1+n\cdot n'\) exactly. Let \(T(\eta)\) be the resulting field
coefficient tensor with prescribed mark indices \(\eta\). Then
\begin{equation}\label{ht:tensorrecipe}
 Z=T(0),\qquad C_{ij}=\sum_{a=1}^3 T(\eta_i=\eta_j=a),
\end{equation}
with all other mark indices zero. Multiply the local moment tensors
by their least positive common denominators and multiply each field
by the positive denominator of its radius; let \(D_0\) be the product
of these denominators. This gives integer coefficient tensors
\(a=3D_0\htcur\), \(b=5D_0\htnext\).

For a separately affine polynomial \(G(u)=\sum_\alpha g_\alpha
\prod_j u_j^{\alpha_j}\), define
\[
 \mathcal O(G)=\sum_\alpha g_\alpha\bigotimes_j\sigma_{\alpha_j},
 \quad \sigma_0=I,\quad
 \sigma_1=\begin{pmatrix}0&1\\1&0\end{pmatrix},\quad
 \sigma_2=\begin{pmatrix}0&-i\\i&0\end{pmatrix},\quad
 \sigma_3=\begin{pmatrix}1&0\\0&-1\end{pmatrix}.
\]
An inversion \(\mathcal I\) negates all coordinates of its listed
field directions. The two unit fields at any unmarked U corner form
one group and must be inverted together. All other fields, including
mean encodings, form single groups.

Use the symmetric pair basis
\(e_+e_+, e_+e_-+e_-e_+, e_-e_-\), without normalization, for each
unit-field pair; use \(e_+,e_-\) for singles. Let \(V\) be their tensor
product. Their doubled weights are \((2,0,-2)\) and \((1,-1)\),
respectively. On the subspace of total weight \(\kappa\), where
\(\kappa=0\) for an even number of singles and \(\kappa=1/2\) for an
odd number, take the Hermitian-form matrices \(A,B\) of
\(\mathcal O(a^{\mathcal I}),\mathcal O(b^{\mathcal I})\).
For \(c=p/q\), reconstruct
\begin{equation}\label{ht:matrix}
 M=5qA-3pB=15qD_0\,
       \left.\mathcal O(F_c^{\mathcal I})\right|_{\rm central\ form}.
\end{equation}
The equality specifies forms in the displayed integer-coordinate basis,
not matrices in an orthonormal basis.
\end{htdef}

\begin{htlem}[The 16 matrix certificates]\label{ht:matrixcert}
For the 16 non-scalar rows the inversion, dimension, normalization and
Gram parameters are those of Table~\ref{ht:certtable}. Reconstruct
\(M\) by Definition~\ref{ht:reconstruct}. The finite assertions are
\begin{equation}\label{ht:gramchecks}
 M=gN,\qquad
 E=d\,s^2N-LL^{\mathsf T},\qquad
 E_{ii}>\sum_{j\ne i}|E_{ij}|\quad\hbox{for every row }i,
\end{equation}
where \(g,d,s\) are positive integers, \(N\) is an integer symmetric
matrix, and \(L\) is an integer lower triangular matrix. Here every
\(d=1\). The supplied \(N,L\) arrays are in
\texttt{data/\allowbreak universal\_certificates.jsonl}. The equality \(M=gN\) is
checked entry by entry against the reconstructed matrix; the arrays
are not assumed to have the correct mathematical origin.
These assertions imply all 16 non-scalar inequalities of H1-tree.
\end{htlem}
\begin{proof}
First we justify the reconstruction and its use of a central form.
At a marked U vertex, path degree plus field count plus one insertion
is at most four; at an unmarked U corner it is four. At the unmarked
D vertex \(S\) it is two. Mark insertions occur at two distinct vertices,
never twice at one vertex. Thus \eqref{ht:moments} and
\eqref{ht:tensorrecipe} give the exact coefficients, with a positive
scale \(D_0\). They give 1,465,600 dense tensor cells and 122,884
nonzero cells, counting both tensors and all 38 rows.

The proof of P.Lemma D.7 applies to these forest functionals: separate
degree at most four gives orthogonal covariance; conjugating all Pauli
matrices by the same element of \(\mathrm{SU}(2)\) therefore leaves
the operator unchanged. It commutes with the total spin operators
\(S_z,S_+,S_-\). Reflection of all second coordinates gives a real
symmetric matrix. At a corner with two unit fields, the local Pauli
factor is \((I\pm n\cdot\sigma)^{\otimes2}\). The spin \(n\) is
unit, so this factor has range in \(\operatorname{Sym}^2\mathbb C^2\)
and vanishes on its orthogonal complement. Inverting the entire pair
preserves this property. Mean encoding creates no pair. There are no
retained D field pairs in these models. Consequently the supported
space is precisely \(V\), with just one central form per model.
P.Lemma D.9, for the concrete single-spin and symmetric-pair generators,
promotes positivity on the weight-\(\kappa\) space to positivity on all
of \(V\). These arguments concern the functionals themselves and do
not assume that a forest came from a covered pattern.

The checker \texttt{checks/reconstruct.py} constructs the integer tensors
from \eqref{ht:moments}, contracts bonds as in \eqref{ht:tensorrecipe},
forms the Pauli matrices, and sandwiches each pair with the displayed
integer basis. The exact replay also checks real symmetry and the
vanishing of entries between different weights. It compares every
entry in \eqref{ht:matrix} to \(gN\), checks the triangular shape of
\(L\), and tests every residual row in \eqref{ht:gramchecks}.
There are 14,502 matrix cells, 7,443 triangular Gram entries,
129,491 triangular Gram products, and 384 residual rows. Normalized
matrix entries require at most 37 bits and Gram entries at most 19 bits.
No floating point factorization is used in this replay.

For completeness, for real \(x\), row dominance gives
\[
 x^{\mathsf T}Ex\ge
 \sum_i\left(E_{ii}-\sum_{j\ne i}|E_{ij}|\right)x_i^2>0
 \quad(x\ne0),
\]
using \(2|x_ix_j|\le x_i^2+x_j^2\). Hence
\(d s^2x^{\mathsf T}Nx=|L^{\mathsf T}x|^2+x^{\mathsf T}Ex>0\).
Positive scaling gives \(M>0\), also for complex forms by splitting
real and imaginary parts. Equation~\eqref{ht:matrix} and the preceding
central-weight argument give
\(\mathcal O(F_c^{\mathcal I})\ge0\) on the full tensor space,
with zero action on discarded antisymmetric pair spaces.
Finally P.Lemma D.6, or directly the Pauli trace identities, gives
\[
 F_c(u)=\operatorname{Tr}\left(
 \mathcal O(F_c^{\mathcal I})
 \bigotimes_j\frac{I+(u^{\mathcal I})_j\cdot\sigma}{2}\right)\ge0
 \qquad (|u_j|\le1).
\]
The factors are positive semidefinite probability densities. This proves
the inequality on the required balls.
\end{proof}
\begin{longtable}{@{}rlrrrrr@{}}
\caption{All central forms. Inversions use zero-based field indices in Definition~\ref{ht:reconstruct}. The margin column is the minimum residual row margin; every row is checked.}\label{ht:certtable}\\
\toprule Model & \(\mathcal I\) & Dimension & \(g\)&\(d\)&\(s\)&Margin\\\midrule\endfirsthead
\toprule Model & \(\mathcal I\) & Dimension & \(g\)&\(d\)&\(s\)&Margin\\\midrule\endhead
\bottomrule\endfoot
5 & \(0,1,2,3\) & 14 & 500 & 1 & 1 & 2056788\\
12 & \(1,2,3,4\) & 7 & 300 & 1 & 1 & 250381\\
13 & \(0,1,2,3\) & 25 & 600 & 1 & 1 & 991054\\
14 & \(2,3,4,5\) & 14 & 1500 & 1 & 1 & 26056\\
15 & \(4,5\) & 14 & 8 & 1 & 1 & 133994474\\
16 & \(1,2,3,4,5,6\) & 25 & 120 & 1 & 1 & 8293235\\
17 & \(1,2,3,4\) & 7 & 3000 & 1 & 1 & 4030\\
18 & \(1,2,3,4,5\) & 14 & 240 & 1 & 16 & 1847052\\
19 & \(4,5,6,8\) & 66 & 400 & 1 & 1 & 30237860\\
20 & \(2,3,4,5\) & 36 & 1000 & 1 & 1 & 194208\\
26 & \(0,1\) & 14 & 2500 & 1 & 1 & 2062552\\
33 & \(4,5\) & 25 & 600 & 1 & 1 & 942356\\
34 & \(1,2,3\) & 14 & 60 & 1 & 1 & 2037189\\
35 & \(1,2,3,4\) & 7 & 60 & 1 & 1 & 977583\\
36 & \(0,1,7,8\) & 66 & 400 & 1 & 1 & 1068222876\\
37 & \(0,1,2,3\) & 36 & 16 & 1 & 1 & 133061543\\
\end{longtable}

\begin{htprop}[Discharge of the finite tree input]\label{ht:certified}
The exact rational assertions in Lemmas~\ref{ht:scalar} and
\ref{ht:matrixcert} certify H1-tree. Lemmas~\ref{ht:meanbern} and
\ref{ht:specialmean} supply the additional finite mean inequalities
used in reducing actual passages to that input.
\end{htprop}
\begin{proof}
The 22 scalar rows and 16 non-scalar rows partition Table~\ref{ht:modeltable}.
The two lemmas prove the required inequalities for every direction,
not just at sampled fields. The recorded exact replays verify all
assertions stated in these lemmas. The mean lemmas are separate inputs
to the reduction, not consequences of matrix positivity.
\end{proof}

\begin{htdef}[Actual local data]\label{ht:actual}
Fix a finite \(G\subset\mathbb Z^2\) and original leg pins as in
P.Definition 2.3, with protected set
\(S(G)=\{x\in G:d(x,\mathbb Z^2\setminus G)\le2\}\).
Consider a U departure at \(y\notin S(G)\), with its incoming bond
already absent. Normalize by a lattice symmetry so that \(y=0\),
the arrival site is \(S=(0,-1)\), and the outgoing site is
\(o\in\{E,N,W\}\). The non-exceptional local hypothesis is precisely
the non-tip kept-set hypothesis of P.Definition 6.8:
\begin{enumerate}
\item The passage is U, keeps its three non-arrival bonds, and its
chosen outgoing bond is present.
\item The arrival site is removed, or is D keeping two of its three
non-passage bonds.
\item Every other nonremoved neighbour is U with full kept set, or
D keeping its passage bond and one other bond.
\item Every other nonremoved non-tip U site has full kept set; every
D corner keeps two bonds. Between nonremoved non-tip ends, kept
flags agree.
\end{enumerate}
Tips have their actual ordinary-domain statuses and kept sets; a missing
tip bond is allowed. All original legs are retained, and D sites carry
none. There are no artificial corner slots in these data.
Write \(E_{\rm net}\) for the network before the step and
\(E_{\rm net}^{+}\) for the network after changing the passage to D
and deleting its outgoing bond, as in P.Definition 6.2.
\end{htdef}

\begin{htlem}[Pin and field counts at an actual passage]\label{ht:counts}
The radius-two template lies in \(G\). The passage and its four nearest
neighbours carry no original leg. Each corner has at most two external
factors, each actual tip at most three. A D corner with \(k\) present
internal bonds has at most \(2-k\) external factors. The same bounds
hold after fixing arbitrary outside spins of norm at most one, and
all resulting fields have norm at most one.
\end{htlem}
\begin{proof}
Since \(y\notin S(G)\), its distance from the complement is at least
three. Its neighbours have distance at least two, whereas an original
leg can occur only at distance one. Thus no original leg is incident
to the passage or its neighbours. Each corner has two lattice directions
inside the template and two outside; each tip has one inside and three
outside. In any one direction there is at most one factor: an original
leg to the complement, or a present bond to a site of \(G\), or no factor.
Deleting sites creates no new pins (P.Definition 2.3).

Original legs occur only at protected U sites by P.Definition 6.1.
The trace never departs from a protected site, by P.Definition 6.2,
so it cannot create a D site carrying a pin. A D corner therefore has
only present-bond factors. Its two-element kept set bounds the number
of its present bonds by two; subtracting its \(k\) internal bonds proves
the sharper external count. Kept bonds to removed sites contribute no
factor and cannot invalidate this upper bound. Outside spins have norm
at most one by the supports of both laws; pins are unit vectors.
Fixing them produces the asserted field factors.
\end{proof}

\begin{htlem}[The reduction identities and all local cases]
\label{ht:localcases}
After fixing the outside spins, an actual current/next pair is a
nonnegative combination, with identical coefficients for its two
members, of positive common scalars times tree-model current/next pairs.
The tree models before radius enlargement are constructed by the
following exhaustive procedure.
\end{htlem}
\begin{proof}
All sums are finite, so their order can be changed. Denote the set of
present non-arrival neighbours by \(M_0\subset\{E,N,W\}\), including
\(o\). Integrate the passage first. It has no pin and at most three
bonds. Equation~\eqref{ht:moments} gives \eqref{ht:polynomials} with
\(M=M_0\), exactly as in P.Definition 6.5 and P.Lemma D.1. The next
expression has the same base functional and insertion only between
the two non-outgoing neighbours, when both exist. In particular none
of the following base-measure transformations depends on current
versus next.

\emph{Step 1: tips.}
For each nonremoved actual tip, integrate its external weight \(W_t\)
first. It has at most three factors, with radii at most one. Put
\(A_t=\lambda_{\mathrm{st}(t)}(W_t)>0\) and
\(b_t=\lambda_{\mathrm{st}(t)}(nW_t)/A_t\).
Lemma~\ref{ht:elementary} gives \(|b_t|\le3/5\), also for a D tip.
Its actual internal bond, if present, becomes \(1+b_t\cdot n_x\),
with common scalar \(A_t\). If the bond is absent there is only that
scalar. A removed tip contributes one. If a non-tip kept flag points
to a missing or nonbonded actual tip, assign its prospective field zero.
This is permitted by affinity and preserves the exact current/next
expressions. Thus all required covering tip fields may be given radius
\(3/5\), without constructing fictitious pins or changing a D kept set.
The arrival tip, when bonded, also gives a field of radius \(3/5\).
In all subsequent combinatorial references, a tip bond is called present
when the non-tip kept flag requires it, even if its effective field is
zero because the actual tip bond was absent. Non-tip bonds retain their
actual meaning. This convention is exactly the covering-tip convention
of P.Definition 6.8.

\emph{Step 2: D passage neighbours and upper D corners.}
A D neighbour has only its passage bond and at most one additional
present bond. If the latter goes to its tip, after Step 1 it is an
isolated marked D site with one field of radius at most \(3/5\).
Keep its D law; its mean has radius at most \(3/25\).
For every other D neighbour expand its law as
\begin{equation}\label{ht:dlaw}
 \lambda_{\htD}=\tfrac35\lambda_{\htU}+\tfrac25\delta_0.
\end{equation}
In the zero term every incident bond is 1, and every passage insertion
at that neighbour is zero. Remove that variable and its mark; remove
the outgoing mark as well if this was \(o\). In the U term retain the
variable with its original other bonds. Let \(M\) denote the surviving
marked set. In particular a lost outgoing mark does not turn another
mark into an outgoing mark.

An upper corner \(A\) or \(B\) with two original internal bonds and
law D has no external factor, by Lemma~\ref{ht:counts}. If its two
neighbour spins are \(a,b\), its integrated kernel is
\begin{equation}\label{ht:topd}
 1+\tfrac15a\cdot b
   =\tfrac35\lambda_{\htU}((1+n\cdot a)(1+n\cdot b))
       +\tfrac25.
\end{equation}
Use these two terms with weights \(3/5,2/5\). Equivalently expand the
corner law before its integration. These choices can be made before or
after \eqref{ht:dlaw}, by finite distributivity. A corner zero term
deletes that bridge. If one endpoint is zero, the active U term has
one bond and no field, and integrates to exactly 1; if both are zero
it does so as well. It must not be replaced at this stage by a U leaf
with two arbitrary fields.

Every resulting term has coefficient
\((3/5)^a(2/5)^b\), the product over these specified binary choices.
It is nonnegative, independent of spins, and the same for current and
next. If \(M=\varnothing\), both polynomials are the same positive
base normalizer and \(F_c=(1-c)Z\ge0\); set aside this term.
The D neighbours kept at their tips are isolated, since they have no
other kept bond after passage deletion. Integrating each gives the
common normalizer 1 and its marked mean of radius \(3/25\).

\emph{Step 3: the lower part and the arrival site.}
Write \(L=(-1,-1)\) and \(R=(1,-1)\). The lower graph is a subgraph
of \(W-L-S-R-E\), with a possible arrival-tip field at \(S\).
There are no other arrival factors and no arrival pin. The following
cases exhaust its possibilities, because S is removed or D with two
kept bonds.

(a) If both bonds \(SL,SR\) are present, the arrival tip bond is absent.
Integrating S gives
\begin{equation}\label{ht:arrivalcut}
 1+\tfrac15n_L\cdot n_R
   =\frac1{12}\sum_{\xi\in D}
        (1+\sqrt{3/5}\,\xi\cdot n_L)
        (1+\sqrt{3/5}\,\xi\cdot n_R).
\end{equation}
This identity follows from the first and second moments of \(\xi\).
It is a positive representation: each factor is nonnegative since
\(\sqrt{3/5}<1\). The U law used here is an auxiliary finite summation,
not an independence assumption about the old corner spins.

If both E and W survive, use \eqref{ht:arrivalcut}, except in case
(d) below. On each fixed \(\xi\) fibre, a U lower corner attached to
its surviving side mark has its at most two external fields, each of
radius one, and this one auxiliary field of radius \(\sqrt{3/5}\).
It has at most three fields before its bond to the side mark is
integrated, so its mean has norm at most \(3/5\), by P.Lemma 2.5.
A D lower corner attached to its side mark has the two internal bonds
to that mark and to S in the original configuration. It has no external
field; after the split its only field is the auxiliary one. Its mean
is exactly \(\sqrt{3/5}\,\xi/5\), of norm at most \(1/5\).
Integrate each corner by \eqref{ht:leaf}. A corner not attached to a
surviving mark contributes just its strictly positive normalizer.
Thus the two effective side fields are conditionally separate. The
coefficient for a fixed \(\xi\), after both corner integrations, is
\begin{equation}\label{ht:splitweights}
 \frac1{12}A_L(\xi)A_R(\xi)>0.
\end{equation}
Neither normalizer includes a side-mark factor or a passage insertion.
They depend only on the fixed outside fields and \(\xi\), so these
are common coefficients for current and next. One need not normalize
these coefficients to sum to one; dividing by their common positive
sum would give the same convex statement.

(b) Suppose both \(SL,SR\) are present, but only one of E,W survives.
For example let E survive. First integrate the opposite corner L,
which has no bond to a surviving side mark. A U L has at most its two
unit-ball external fields, hence mean norm at most \(1/2\). A D L has
its bond to S and therefore at most one external field, hence mean norm
at most \(1/5\le1/2\). (A bond formerly incident to a zero mark is now
1 and adds no field.) With \(\mu_L\) this mean, the subsequent D
arrival integral is exactly
\begin{equation}\label{ht:onealive}
 \lambda_{\htD}((1+\mu_L\cdot n_S)(1+n_S\cdot n_R))
       =1+(\mu_L/5)\cdot n_R.
\end{equation}
The effective arrival field at R thus has radius at most \(1/10\).
If R is U and attached to E, it has its two external unit-ball fields
and this small field, so Lemma~\ref{ht:specialmean} gives radius
\(18/35\) for the field it sends to E. If R is D and attached to E,
its original two internal bonds leave no external field, and its sole
new field has radius at most \(1/10\); its mean has radius at most
\(1/50\le1/5\). If it is not attached to E, discard its common
positive normalizer. The case with W surviving is the reflected
calculation, integrating R, then S, then L. If neither side survives,
the entire lower component contributes a common positive scalar.

(c) Suppose the two-lower-bond case does not hold. If S has one lower
bond and its tip bond, first integrate the tip as in Step 1, then S:
\begin{equation}\label{ht:arrivaltip}
 \lambda_{\htD}((1+b_S\cdot n_S)(1+n_S\cdot n_C))
       =1+(b_S/5)\cdot n_C,\qquad |b_S/5|\le3/25.
\end{equation}
Here \(C\) is the attached lower corner. If it is U and attached to
its surviving side mark, its fields have radii at most \((1,1,3/25)\),
so its mean has radius at most \(137/265\) by
Lemma~\ref{ht:specialmean}. If it is D and attached to that mark,
its two original internal bonds leave no external factor; its new
mean is at most \(3/125\le1/5\).
If S has a lower bond but no present tip bond and no other lower bond,
its integration is simply \(\lambda_{\htD}(1+n_S\cdot n_C)=1\).
A removed S has the same effect. A remaining attached U lower leaf
then has at most two external unit-ball fields, giving radius \(1/2\).
An attached D lower leaf has at most one external field when its
arrival bond was absent, or no external field when that bond was
present. Its mean radius is at most \(1/5\). If S has no lower bond,
its component is a common positive scalar (with only its possible
one tip field its normalizer is 1). Lower components not connected
to any surviving mark are likewise discarded. These statements also
cover kept arrival bonds whose other endpoint was removed.

(d) The correlation-preserving case is exactly
\begin{equation}\label{ht:preserve}
 E,W\in M,\quad N\notin M,\quad SL,SR\text{ present},\quad
 L,R\text{ have original law U}.
\end{equation}
Do not integrate S, L or R; retain
\(E-R-S-L-W\), with the D law at S and two original corner fields
at each of L,R. Consistency forces the side bonds \(ER,LW\):
a surviving E or W is nonremoved before the zero-law expansion, and
a U corner keeps its side bond, so the other nonremoved endpoint
must keep it too. A surviving U term of an originally D endpoint
inherits that present bond. The arrival has no tip field because it
keeps the two lower bonds. The retained five-vertex path is therefore
exact, with no split and no discarded correlation. If N survives,
case (a) is used even when both lower corners are U. If one lower
corner is D, case (a) is used when E,W survive and case (b) otherwise.
These are all alternatives to \eqref{ht:preserve}.

The need to retain this correlation can be seen without optimization.
If one cuts both lower fields to independent balls of radius \(3/5\)
in a turning configuration with E and W surviving and N absent, allowing
also each tip radius \(3/5\) and each upper-leaf radius \(1/2\),
the last row of Table~\ref{ht:meantable} gives two independent mean
balls of radius \(q=226/495\). Their admissible assignment
\(m_E=qe_1,\ m_W=-qe_1\), with E outgoing, has
\[
 \htcur/\htnext=1-q^2/3=683999/735075<187/200.
\]
This is a counterexample to that enlarged independent-mean inequality,
not to the actual local inequality: the enlargement discarded the lower
path constraint. Keeping the path in \eqref{ht:preserve} avoids precisely
this false intermediate assertion.

For reference, the complete side-field rule, whenever a lower corner
is attached to a surviving side mark and is not retained in (d), is
\begin{equation}\label{ht:lowerrule}
 r_{\rm lower}=\begin{cases}
 1/5,&\text{original corner D},\\
 3/5,&\text{original corner U, }SL,SR\text{ present, both side marks survive},\\
 18/35,&\text{original corner U, }SL,SR\text{ present, other side mark absent},\\
 137/265,&\text{original corner U, its arrival bond and arrival tip bond present},\\
 1/2,&\text{original corner U, all remaining cases}.
 \end{cases}
\end{equation}
No side bond means no such field.

\emph{Step 4: the upper part.}
Consider each of A,B separately, after Step 2. An original U corner
has at most two external unit-ball fields. If both its side marks
survive, retain the U bridge and its two fields, inserting zeros for
missing fields. If exactly one survives, integrate this leaf and use
radius \(1/2\); if none survives, discard its positive normalizer.
An original D corner with two internal bonds was already treated by
\eqref{ht:topd}: retain its active U bridge only when both endpoints
survive. It then has no external fields; insert two zero fields.
If fewer survive it gives exactly 1 as established in Step 2.
An original D corner with one internal bond has at most one external
field, by Lemma~\ref{ht:counts}. If its endpoint survives, integrate
it and use radius \(1/5\); otherwise discard its positive normalizer.
A corner with no original internal bond is an isolated common positive
factor. This exhausts original laws U,D,X and all upper internal degrees
0,1,2. Only the retained bridges will subsequently have their two
field radii enlarged to \((1,1)\).

\emph{Step 5: path components and isolated marks.}
Outside \eqref{ht:preserve} the only remaining bridges are subpaths of
\(E-A-N-B-W\). Every such component is a path, with marked endpoints
and possibly marked N in its interior. In \eqref{ht:preserve}, N is
absent, so neither upper corner can join two marks. After its leaf
integration the only nontrivial component is the retained lower path.
Thus no cycle remains. The fixed D tip-neighbours are isolated as
already noted.

At a remaining marked U vertex let \(d_v\) be its path degree. Before
zero padding it has at most \(3-d_v\) fields. Indeed an original U
neighbour has at most its three non-passage directions, each supplying
either a path bond or one integrated field. An activated D neighbour
has only its one non-passage kept direction and satisfies the stronger
bound with one in place of three. All fields are among the tip,
lower-corner, and upper-leaf fields just specified. If \(d_v=0\),
its at most three fields have one of the U radius triples in
Table~\ref{ht:meantable}, after zeros are inserted. Integrate the site
into a marked mean vertex, multiply by its strictly positive normalizer,
and use Lemma~\ref{ht:meanbern}. The assertion about the finite list
of triples is verified explicitly in Lemma~\ref{ht:coverage} below.
The fixed D neighbour uses the D row \((0,0,3/5)\) of that table.

If \(d_v\ge1\), pad its field list to \(3-d_v\) by zero fields.
An unmarked retained U corner has two path bonds and is padded to two
fields. The unmarked retained D arrival has two path bonds and no
fields. These bounds keep the separate degree, including one mark
insertion, at most four for U sites and at most two for D sites.
All normalizers removed in Steps 1--5 are positive and common.
Every binary law expansion has fixed weights \(3/5,2/5\); every use
of \eqref{ht:convexfield} uses coefficients independent of current
versus next; every auxiliary-spin coefficient is
\eqref{ht:splitweights}. Hence all combinations have identical
nonnegative coefficients for the two polynomials, as asserted.
\end{proof}

\begin{htlem}[Finite coverage directly from the kept-set definition]
\label{ht:coverage}
Apply the procedure of Lemma~\ref{ht:localcases} to all actual local
data of Definition~\ref{ht:actual}, using a prospective tip field as
specified in Step 1. Every nonempty reduced forest, up to relabeling
of components and reversal of paths preserving laws and marks, has
field radii componentwise no larger than one row of
Table~\ref{ht:modeltable}. There are 20 decorated geometries and
38 maximal radius lists. The isolated-mean calls are exactly the 21
rows in Table~\ref{ht:meantable}. This is a finite combinatorial
assertion with an exact exhaustive check from the following rules,
not a comparison with the archived pattern catalogue.
\end{htlem}
\begin{proof}
Here are full index sets and a reproducible finite verification of the
remaining radius bookkeeping. They also give a direct method to check
this small-case reduction independently.

Encode a kept set by its mask in direction order E,N,W,S. A corner has
options U with mask 15, X with mask 0, or D with one of the six
cardinality-two masks. The arrival has X or D keeping any two of its
three directions other than its passage direction. Each other neighbour
has U, X, or D keeping the passage direction and one of its other three
directions. Thus there are exactly
\(4\cdot5^3\cdot8^4=2,048,000\) complete raw choices. For a nonremoved
corner and a nonremoved nearest neighbour retain a choice exactly when
the two kept flags agree. These are all consistency comparisons: the
remaining non-tip bonds join the passage to its neighbours, whose flags
are already prescribed. A corner or neighbour that is removed imposes
no comparison. Put in every bond with two nonremoved keeping ends and
every required covering tip bond. Mark every possible outgoing site.
Reflecting \((x_1,x_2)\mapsto(-x_1,x_2)\) reduces the outgoing choices
to N and E, without changing straight/turning type or any field bound.

Delete isolated unmarked factors. Their positive normalizers are common
by Lemma~\ref{ht:elementary}. The remaining effective record consists
only of the present bonds, laws at their non-tip ends, and outgoing
site. Distinct original kept masks with that record have the same
maximal field allowances: two at U corners and \(2-k\) at D corners,
where \(k\) is their present internal degree. The missing actual
factors are zero fields in those allowances. Thus collapsing such
records cannot omit any relevant factor.

For clarity, the entire calculation on each retained record is the
following finite recipe; every operation in it was justified in
Lemma~\ref{ht:localcases}.
\begin{enumerate}
\item Select all binary choices in \eqref{ht:dlaw} for D marks except
those keeping their tips, and in \eqref{ht:topd} for the two-bond upper
D corners. Record the surviving marked set \(M\). The empty set is
trivial and produces no model.
\item At each surviving marked vertex append a radius \(3/5\) if its
tip flag is kept. If it has its lower-corner bond, append the radius
in \eqref{ht:lowerrule}, unless \eqref{ht:preserve} holds; in that case
retain the lower path instead. For an upper corner of original law U,
retain a U bridge when both endpoints survive, append \(1/2\) to its
single surviving endpoint otherwise, and do nothing when neither
survives. For an upper corner of original law D with two original
internal bonds, retain the bridge only in its active binary term and
only with both endpoints surviving. For an original one-bond D upper
corner append \(1/5\) at its surviving endpoint. There are no other
bridges or fields.
\item Put radii \((1,1)\) at every retained unmarked U corner;
this enlarges the zero fields of an activated D corner if necessary.
Retain law D and no fields at S in the lower path. A nontrivial path's
marked vertices are U: a D mark kept at its tip is isolated. Pad a U
mark of path degree \(d_v\) to \(3-d_v\) fields with zeros. For each
isolated mark, pad to three fields and replace its list by its mean
radius \eqref{ht:meanformula}. Record every such call. The only D
call is \((0,0,3/5)\).
\item Sort the radii at each vertex. A vertex attribute is
\((\text{law},\text{marked},\text{outgoing},\text{radius list})\),
where a mean has law-symbol m and a single radius. For each path take
the lexicographically smaller of its two orientations, comparing first
its decorated geometry (the first three attributes and field counts)
and then its radius lists. Sort components the same way. Keep the
straight/turning flag as part of the key.
\item Within a fixed key compare corresponding radii componentwise.
Discard a list dominated by another, retaining each maximal list once.
Verify that the results are exactly Table~\ref{ht:modeltable} and that
every mean call belongs to Table~\ref{ht:meantable}.
\end{enumerate}
For a corner/mark with a smaller actual external field, the appropriate
step uses its ball bound rather than an assumption of maximal norm.
Thus this recipe is an upper-radius computation, not a claim that all
choices of enlarged fields are jointly realizable by a global domain.
The inequality is being proved on the larger independent product of
balls, which covers those correlated choices as well.

The exact enumeration and comparisons are implemented without catalogue
imports in \texttt{checks/\allowbreak audit.py}, function \texttt{coverage}.
It filters the above mask choices, computes the effective records from
these choices, and carries out Steps 1--5 with rational radii. It also
checks every generated list against a dominating row and verifies
both inclusions of the maximal-list equality, not just the count 38.
Its complete counts are
\begin{equation}\label{ht:coveragecounts}
 \begin{array}{lr}
 \text{consistent raw choices}&106101\\
 \text{choices with any marked outgoing site}&184914\\
 \text{choices with outgoing N or E}&123339\\
 \text{effective reflection classes generated from those choices}&7962\\
 \text{binary reduction branches on those classes}&24741\\
 \text{branches with no surviving mark}&981\\
 \text{decorated geometries}&20\\
 \text{maximal radius lists}&38.
 \end{array}
\end{equation}
The coincidence of the effective-class count with the old catalogue's
size is a check of the definitions, not an input. No entry or identifier
from that catalogue is read. The geometries are exhausted by the
following partition of the maximal rows; a path has a U corner between
consecutive displayed marks, except for the retained lower path.
\begin{center}\small
\begin{tabular}{cll}
\toprule Type & Geometry and outgoing location & Maximal rows\\\midrule
 t & one mean, outgoing present & 0\\
 t & two means, outgoing present & 1,2,3\\
 t & one mean, outgoing absent & 4\\
 t & three-vertex path, outgoing at an end & 5\\
 t & three means & 6,7,8\\
 t & two means, outgoing absent & 9,10,11\\
 t & mean plus three-vertex path, outgoing on path & 12,13,14\\
 t & three-vertex path, outgoing absent & 15\\
 t & outgoing mean plus three-vertex path & 16,17,18\\
 t & upper five-vertex path, outgoing at an end & 19\\
 t & lower five-vertex path, outgoing at an end & 20\\
 s & one mean, outgoing absent & 21\\
 s & two means, outgoing present & 22,23,24\\
 s & one mean, outgoing present & 25\\
 s & three-vertex path, outgoing at an end & 26\\
 s & two means, outgoing absent & 27,28\\
 s & three means & 29,30,31,32\\
 s & mean plus three-vertex path, outgoing on path & 33,34,35\\
 s & upper five-vertex path, outgoing in the middle & 36\\
 s & lower five-vertex path, outgoing absent & 37\\\bottomrule
\end{tabular}
\end{center}
This states exactly the finite coverage assertion used in the proof,
including its computer-checked part. The check uses only finite products,
Boolean mask comparisons, rational arithmetic and list comparisons.
It introduces no numerical inequality or catalogue-exhaustiveness
assumption.
\end{proof}

\begin{htlem}[Relabeling and radius extension]\label{ht:extension}
A relabeling of path vertices preserving laws, marks and the outgoing
mark preserves \(Z,C_{ij},\htcur,\htnext\), after the corresponding
field relabeling. Componentwise radius domination as in
Lemma~\ref{ht:coverage} transfers H1-tree to the dominated model.
\end{htlem}
\begin{proof}
Relabeling is a bijection on the finite product variables and preserves
every scalar product and site law. Reversing a path and permuting
components are particular examples. The type flag is retained
separately. Fields at the same vertex can be permuted because their
factors commute.

For a smaller field vector in a larger listed ball, use
\eqref{ht:convexfield}. The polynomial \(F_c\) is separately affine
in each field vector. At an isolated mean vertex it is separately
affine in that mean as well, because every passage pair insertion
uses a marked vertex only once. Expanding each vector into a convex
combination of boundary vectors therefore expresses the smaller value
of \(F_c\) as a convex combination of values on the larger model.
These are nonnegative by H1-tree. A new zero field contributes the
factor 1 exactly and can be treated by the same identity. No division
by a partition function or ratio monotonicity is involved.
\end{proof}

\begin{htthm}[Reduction to the actual conclusion of (H1)]
\label{ht:reduction}
Assume H1-tree. Use the mean inequalities of
Lemmas~\ref{ht:meanbern} and \ref{ht:specialmean} and the finite
coverage assertion of Lemma~\ref{ht:coverage}, established by the
specified exact checks above. For every actual non-exceptional U departure
of Definition~\ref{ht:actual}, every original unit pin assignment
\(\tau\), and every fixed outside configuration of spins of norm
at most one, the two conditional template sums satisfy
\begin{equation}\label{ht:conditional}
 C(\omega)\ge c_{\rm type}N(\omega).
\end{equation}
Consequently, for the full networks,
\begin{equation}\label{ht:global}
 Z_{E_{\rm net}}^\tau\ge c_{\rm type}
                         Z_{E_{\rm net}^{+}}^\tau,
 \qquad \rho_j^\tau\le1/c_{\rm type}.
\end{equation}
The finite checks specified above discharge all these assumptions.
In particular this proves exactly the local bound used in Section 6
with the constants \eqref{ht:constants} unchanged.
\end{htthm}
\begin{proof}
The template lies in G by Lemma~\ref{ht:counts}. Fix all spins outside
it. Put every site weight and factor wholly outside it into
\(H(\omega)\ge0\). These outside factors are identical before and
after the step, since the only changes are at the passage and its
outgoing bond, both inside the template. The remaining two template
sums are \(C(\omega),N(\omega)\). Their external fields obey precisely
the counts in Lemma~\ref{ht:counts}, irrespective of correlations
between the fixed outside spins or coincidences among pin directions.

Lemma~\ref{ht:localcases} integrates and decomposes this pair. Term
by term its difference is a nonnegative coefficient times a positive
common scalar times the difference of one of its reduced forests.
Empty marked sets have nonnegative difference \((1-c)Z\). For every
other forest Lemma~\ref{ht:coverage} supplies a row of the 38-model
list with dominating radii, and Lemma~\ref{ht:extension} transfers that
row's inequality to this forest. Thus every summand is nonnegative.
This proves \eqref{ht:conditional} for every outside assignment,
including those of zero global marginal weight. The positive scalars
need not be uniform in \(\omega\); they are identical for the two
members of each pair, which is the property required here.
Multiplication by \(H(\omega)\), followed by the finite sum over outside
spins, proves \eqref{ht:global}. Both partition functions are strictly
positive by P.Lemma 2.4, so the ratio assertion follows.
Proposition~\ref{ht:certified} and the exact coverage check discharge
the hypotheses without an optimization or sampling assumption.
\end{proof}

\begin{htprop}[How the theorem enters the paper]\label{ht:use}
In the first trace of P.Lemma 7.2 every U departure satisfies the local
hypothesis of Theorem~\ref{ht:reduction}. In the second trace of
P.Lemma 7.3 the same holds outside
\(\{y':0<d(y',e)\le2\}\), where \(e\in S(G)\) is the fixed stopped
arrival of the first trace. There are at most twelve exceptional U
departures, and their existing \(12/5\) ratio bound is unchanged.
Thus all uses of the local input in P.Lemma 6.13 and Section 7 remain
valid on replacing the hypothesis (H1) by the above finite tree input
and its specified mean and coverage checks.
\end{htprop}
\begin{proof}
The kept-set induction and non-tip local classification are exactly
P.Lemma 7.2. They give a live U kept set of size three, an arrival X
or D kept set of size two, ordinary kept sets elsewhere, and agreement
of kept flags between nonremoved ends. Those are
Definition~\ref{ht:actual}. Actual tip treatment in the present proof
requires only its factor count and norm bound, so no covering global
tip status has been assumed. Pins are retained as in
Lemma~\ref{ht:counts}.

P.Lemma 7.3 gives one fixed protected U defect at e in the second trace.
When \(d(y',e)>2\), it is outside the template, whose non-tip data
satisfy the same rules. No departure occurs at e. The square lattice
has four sites at distance one and eight at distance two from e, and
a site is departed from with law U at most once, since its first U
departure changes it to D. There are therefore at most twelve exceptions.
At each exception P.Lemma 6.6 gives \(\rho_j^\tau\le12/5\).
At D departures the ratio is one by P.Definition 6.5. At all other U
departures Theorem~\ref{ht:reduction} gives the same bound
\(1/c_{\rm type}\) that P.Lemma 6.10 supplied. These are the only local
ratio facts used in P.Lemma 6.13 and the subsequent response recursion.
For the type-dependent response coefficients, P.Definition 6.5 writes
\(Z_E^\tau=Z_{E'}^\tau I^\tau\) and
\(Z_{E^+}^\tau=Z_{E'}^\tau J^\tau\).
Equation~\eqref{ht:global} therefore gives
\(I^\tau\ge c_{\rm type}J^\tau\); with two neighbours it gives
\(I^\tau\ge c_{\rm type}\), since \(J=1\).
These are exactly the local inequalities used in P.Lemma 7.6, so its
sharper type-dependent coefficient bounds are also retained.
There is no claim here about the separate proposed H2 or H3 replacement.
\end{proof}

\begin{htrem}[Repairs and scope of the computational evidence]
\label{ht:repairs}
The following qualifications are necessary for a correct replacement.
First, its theorem is \eqref{ht:global}, not the artificially filled
D-slot version of (H1). Second, the two special mean checks
\eqref{ht:specialradii} are needed in addition to the 21 cube checks.
Third, an upper D corner with two original internal bonds that loses
an endpoint gives exactly 1 after its law expansion; two arbitrary U
fields may be added only to a retained bridge. Fourth, the outgoing
mark can disappear in a zero-law term, but the two surviving non-outgoing
marks still supply \(C_{fg}/5\) in the next polynomial. Fifth, a D
neighbour keeping its tip is integrated as D, rather than replaced by
a stronger U term. Sixth, the lower path in \eqref{ht:preserve} is
retained. The report already flags the last two choices; the proof
specifies their exact scope and the common weights. The source-code
field \(3R\) for a mean is only the formal tensor identity in
Definition~\ref{ht:reconstruct}. The older 19 independent boxes do
not certify the present 22 scalar models.

The exact replay in \texttt{checks/audit.py}, using
\texttt{checks/reconstruct.py}, verifies all finite arithmetic and
coverage assertions above. It does not read \texttt{pext7962.json},
the old operator matrices, or the old inversion choices. Its inputs
are the displayed mathematical definitions and the copied rational
witness files. The checks are native exact rational/integer checks;
no Lean-kernel replay is claimed.

As a separate falsification test, \texttt{checks/numerical.py} used
seed 20261116 on 256 randomly selected effective configurations generated
from the kept-set rules and one source configuration for each maximal
row, with six field assignments per configuration. It tested 1,764
local field assignments and 6,336 binary decomposition branches,
including 522 auxiliary-arrival splits and 24 retained-lower-path
branches. External factors were unit or shorter vectors, with only
\(2-k\) possible fields at D corners, and tip fields had radius at
most \(3/5\). The tip-reduced current/next pair was compared separately with
the law mixture and with the arrival/leaf integration chain. All tested
mean bounds passed; the maximum normalized equality error
\(\max|a-b|/\max(1,\max|a|)\) was less than
\(8.89\cdot10^{-16}\), and the smallest sampled ratio margin above
its type threshold was greater than 0.0195. These floating tests are
diagnostics, not proof inputs, and make no assertion about a sharp
minimum. The exact finite checks and the preceding identities supply
the proof.
\end{htrem}
